\usepackage[margin=0.6in]{geometry}
\usepackage{bm}
\usepackage{amsmath,amsthm,amssymb,amsfonts, fancyhdr, color, comment, environ}
\usepackage{xcolor}
\usepackage{mdframed}
\usepackage[shortlabels]{enumitem}
\usepackage{indentfirst}
\usepackage{hyperref}
\usepackage{graphicx}
\usepackage{tikzsymbols}
\usepackage{tcolorbox}
\usepackage{mathpazo}
\usepackage[object=vectorian]{pgfornament}
\hypersetup{
    colorlinks=true,
    linkcolor=blue,
    filecolor=magenta,      
    urlcolor=blue,
}

\renewcommand{\a}{\texttt{a}}
\renewcommand{\b}{\texttt{b}}
\renewcommand{\c}{\texttt{c}}
\renewcommand{\d}{\texttt{d}}

\newcommand{\0}{\texttt{0}}
\newcommand{\1}{\texttt{1}}
\newcommand{\2}{\texttt{2}}
\newcommand{\h}{\texttt{\#}}

\usepackage{tikz}
\usepackage{float}
\usetikzlibrary{automata, positioning, arrows}
\tikzset{
%->, % makes the edges directed
>=stealth, % makes the arrow heads bold
node distance=3cm, % specifies the minimum distance between two nodes. Change if necessary.
every picture/.style=thick,
every state/.style={thick, fill=gray!10}, % sets the properties for each ’state’ node
initial text=$ $, % sets the text that appears on the start arrow
}

\newcommand{\sectionline}{%
  \noindent
  \begin{center}
  {\color{black}
    \resizebox{0.5\linewidth}{1ex}
    {{%
    {\begin{tikzpicture}
    \node  (C) at (0,0) {};
    \node (D) at (9,0) {};
    \path (C) to [ornament=85] (D);
    \end{tikzpicture}}}}}%
    \end{center}
  }


\pagestyle{fancy}


\newenvironment{problem2}[2][Problem]
    { \begin{mdframed}[backgroundcolor=white] \textbf{#1 #2} \\}
    {  \end{mdframed}}

\newenvironment{problem}[5][Problem]
    { \begin{tcolorbox}[colback=white, title= {#1 #2 (#3): #4 (#5 points)}]}
    {  \end{tcolorbox}}

% Define solution environment
\newenvironment{solution}
    {\textit{Proof:}}
    {}

\renewcommand{\qed}{\quad\qedsymbol}

% prevent line break in inline mode
\binoppenalty=\maxdimen
\relpenalty=\maxdimen

\newcommand{\examinfo}[6]{
\rhead{\Huge{\boxed{\textsf{\textbf{#1}}}}}
%\rhead{\inclraphics[scale=.2]{bracuLogo.jpeg}}
\lhead{#2\\ #3} 
\chead{\textsc{#4 \\Total marks: #5}}
}

% default heading
\def\coursecode{Spring 2024}
\def\coursetitle{Automata and Computability}
\def\examtype{Midterm}
\def\marks{40}
\def\duration{80}

\newcommand{\setexamheading}{
\examinfo{\problemset}{\coursecode}{\coursetitle}{\examtype}{\marks}{\duration}
}

\newenvironment{examinstructions}{
    \begin{tcolorbox}
    \begin{center}
}{
    \end{center}
    \end{tcolorbox}
}

%%%%%%%%%%%%%%%%%%%%%%%%%%%%%%%%%%%%%%%%%%%%%%%%%%%%%%%%%%%%%%%%%%%%%%%%%%%%%%%%%%%%
